% Pursuing Stacks (cotes 134-2 to 134-8): archive cote 134-7, site page 88 (PDF page 89), typescript page 530.
% typescript. Draft compared: ps6.tex:4433-4501.
% AI-transcribed by Claude Opus 5.5 (claude-opus-5-5, adaptive thinking, effort high), Message Batches API, batch ps4 (msgbatch_012Aoc1NuskATPsyBeKSYK7J), 2026-11-12; not yet checked by eye.
\scanpage{134-7/88}
\mspage{530}
functors is compatible with the functors K\hand{$\sim$} , but \inserted[typed]{(even if k=Z)} definitely \uline{not} with \xout{its} \inserted{their}
left adjoints L\hand{$\sim$}H. namely with homology. Despite this fact, we hope \inserted{if k=Z} that
\inserted{the functors} (13) are compatible with the functors S\hand{$\sim$} (assuming that S\hand{$\sim$} can actually
be constructed in the discrete set-up too)\hand{,} \deleted{\hand{if k=Z},} so that spheres are
transformed into spheres. Compatibility with L\hand{$\Pi$}. does'nt have a meaning
strictly speaking, as this functor is not defined on ordinary homotopy
types, however the functors \hand{$\Pi$}$_i$ are defined in both contexts, and the
sections functor should commute to these (for arbitrary k). Thus the
only serious uncompatibility trouble concerns the homology functor LH. ,
\margin{\hand{(H)}} and should not arise however for sphere-like objects S(\hand{$H$},n). \xout{-} I for-%
got to state \inserted{from outset} that the two functors (13) are expected to \xout{\illegible{1}} respect
fibration and cofibration sequence structures, of course.

13.10.

125. During the last eight days, I have been \xout{\unsure{involved}} busy with a
number of things, which left little leisure for mathematical ponderings.
I got a letter from H.J. Baues, who had seen my notes, which induced
him to send me his preprint ``On the homotopy classification problem''
(Chapters I to V + chapter Ext). I spent two evenings looking through
parts of these, where Baues carries as far as possible (namely quite far
indeed) the homotopy formalism\xout{x} \xout{with} in the context of his so-called
``cofibration (or fibration) categories'', using as his leading thread
his ideas on ``obstruction theory''. He makes the point that he tried by
his basic notion to pinpoi\unsure{\hand{/}}nt the weakest axiomatic set-up, sufficient
however \unsure{\hand{f}}or develop\deleted{p}ing all the major familiar (and even some not at all
familiar !) notions, operations and statements of usual homotopy theory.
In his letter, he suggested that maybe in any ``universe'' where homotopy
constructions make a sense, one or the other of his two mutually dual
axiomatic set-ups should be around. Such suggestion was of course quite
interesting for my present reflections, as I do have the hope indeed
that there exists a ``universe'' of schematic homotopy types, which may
be described in terms of the models (namely semisimplicial unipotent
bundles) I have been using so far, or at any rate by closely related
kind of models. More specifically, I do hope for the two kind of opera-%
tions to make sense, namely ``integration'' (including homotopy cofibers\xout{)}
and suspensions) and ``cointegration'' (including homotopy fibers and
loop objects), which in Baues set-up should correspond on the category
of models to a structure of a co\unsure{\hand{f}}ibration category and a fibration cate-%
gory respectively. As the kind of models I have been working with
\margin{88}
